% SPDX-License-Identifier: MIT
% Copyright (c) 2026 Evgeny Baulin
% Optimization Methods, Seminar 03: one-page cheat sheet (English).
% Generated by checks/03. Gradient Descent and Automatic Differentiation.ipynb from web/03/js/content-cheatsheet.js and content/students.json. Do not edit by hand.
% Compile with: latexmk -xelatex -interaction=nonstopmode -halt-on-error Cheatsheet_en.tex
\RequirePackage{fix-cm}
\documentclass[10pt,a4paper]{article}

\usepackage{amsmath,amssymb}
\usepackage{fontspec}
\setmainfont{cmunrm.otf}[BoldFont=cmunbx.otf, ItalicFont=cmunti.otf, BoldItalicFont=cmunbi.otf]
\setsansfont{cmunss.otf}[BoldFont=cmunsx.otf, ItalicFont=cmunsi.otf, BoldItalicFont=cmunso.otf]
\setmonofont{cmuntt.otf}[BoldFont=cmuntb.otf, ItalicFont=cmunit.otf, BoldItalicFont=cmuntx.otf]
\usepackage{polyglossia}
\setdefaultlanguage{english}

\usepackage[a4paper,left=8mm,right=8mm,top=7mm,bottom=11mm,footskip=5.5mm]{geometry}
\usepackage{xcolor}
\usepackage{multicol}
\usepackage{enumitem}
\usepackage{ragged2e}
\usepackage{fancyhdr}
\usepackage[colorlinks=true,linkcolor=hseblue,urlcolor=hseblue]{hyperref}

\definecolor{hseblue}{RGB}{16,45,105}
\definecolor{hseaccent}{RGB}{198,40,40}
\definecolor{hsegray}{RGB}{108,114,126}
\definecolor{hselight}{RGB}{240,243,248}

\hypersetup{pdftitle={Seminar 03 cheat sheet: Gradient Descent and Automatic Differentiation}, pdfauthor={Evgeny Baulin}}

% Page style: the author and the license in the footer.
\pagestyle{fancy}
\fancyhf{}
\renewcommand{\headrulewidth}{0pt}
\fancyfoot[L]{\scriptsize\color{hsegray}\textcopyright\ 2026 Evgeny Baulin, \href{https://t.me/tarakan_tuc}{t.me/tarakan\_tuc}}
\fancyfoot[R]{\scriptsize\color{hsegray}MIT License}

% Compact layout.
\setlength{\parindent}{0pt}
\setlength{\parskip}{0pt}
\setlength{\columnsep}{4.5mm}
\setlength{\columnseprule}{0.3pt}
\renewcommand{\columnseprulecolor}{\color{hselight}}
\setlength{\emergencystretch}{2em}
\setlength{\RaggedRightRightskip}{0pt plus 4em}
\binoppenalty=2000
\relpenalty=500
\setlength{\abovedisplayskip}{2pt}
\setlength{\belowdisplayskip}{2pt}
\setlist[itemize]{leftmargin=7pt, labelsep=3pt, itemsep=1.2pt, topsep=1pt, parsep=0pt, partopsep=0pt,
  label={\color{hseaccent}\scriptsize\textbullet}}

\makeatletter
\newcommand{\cssec}[1]{\par\addvspace{3.5pt}\noindent
  \colorbox{hselight}{\parbox{\dimexpr\linewidth-2\fboxsep}{\RaggedRight\hyphenpenalty=10000\color{hseblue}\bfseries\strut #1}}%
  \par\nobreak\vspace{1pt}\@afterheading}
\makeatother
\clubpenalty=10000
\widowpenalty=10000

\begin{document}
\thispagestyle{fancy}

\noindent
\begin{minipage}[b]{0.6\textwidth}
  \RaggedRight\hyphenpenalty=10000
  {\color{hseaccent}\bfseries\small Seminar 03 cheat sheet}\\[1pt]
  {\color{hseblue}\bfseries\Large Gradient Descent\\ and Automatic Differentiation}
\end{minipage}\hfill
\begin{minipage}[b]{0.39\textwidth}
  \raggedleft\footnotesize
  {\bfseries Optimization Methods}\\
  HSE University, Faculty of Computer Science, 2026--2027\\
  Course lecturer: Andrey Ignatov\\
  Author: Evgeny Baulin, Telegram \href{https://t.me/tarakan_tuc}{@tarakan\_tuc}
\end{minipage}

\vspace{2pt}
{\color{hseblue}\rule{\textwidth}{0.8pt}}
\vspace{-4pt}

\fontsize{9}{11.4}\selectfont
\begin{multicols}{3}
\RaggedRight


% ---- jensen
\cssec{Jensen's inequality}
\begin{itemize}
  \item \textbf{Convex} $f$ on a convex set: $f\big((1-t)x + ty\big) \le (1-t)f(x) + tf(y)$ for all $x, y$ and $t \in [0, 1]$; the chord lies on or above the graph, no derivatives needed. \textbf{Strictly convex}: strict for $x \ne y$, $t \in (0, 1)$. For $x^2$ the gap is $t(1-t)(x - y)^2$.
  \item \textbf{Three equivalent tests} on an open convex set: chord, the definition; tangent $f(y) \ge f(x) + \nabla f(x)^\top (y - x)$ for a differentiable $f$; curvature $\nabla^2 f(x) \succeq 0$ for a twice differentiable $f$. Strict convexity follows from $\nabla^2 f \succ 0$ everywhere, not conversely: $|x|^3$.
  \item \textbf{Weights} $t_i \ge 0$, $\sum_i t_i = 1$: $f\big(\sum_i t_ix_i\big) \le \sum_i t_if(x_i)$; also $f(\mathbb E X) \le \mathbb E f(X)$; reversed for concave functions: $\log \mathbb E X \ge \mathbb E \log X$. \textbf{One pair disproves} convexity: $x^3$ at $x = -2$, $y = 1$, $t = 0.5$ has the graph ($-0.125$) above the chord ($-3.5$).
\end{itemize}

% ---- constants
\cssec{Constants $\mu$, $L$ and $\kappa$}
\begin{itemize}
  \item $\mu$-\textbf{strongly convex}: $f - \frac\mu2\|\cdot\|^2$ is convex, or $\nabla^2 f \succeq \mu I$. $L$-\textbf{smooth}: $\|\nabla f(x) - \nabla f(y)\| \le L\|x - y\|$, or $-LI \preceq \nabla^2 f \preceq LI$.
  \item \textbf{Two parabolas}: $\frac\mu2\|y - x\|^2 \le f(y) - f(x) - \nabla f(x)^\top (y - x) \le \frac L2\|y - x\|^2$. In Jensen form the gap between chord and graph lies between $\frac\mu2 t(1-t)\|x - y\|^2$ and $\frac L2 t(1-t)\|x - y\|^2$.
  \item \textbf{Quadratic} $f = \tfrac{1}{2} x^\top Qx - b^\top x$, $Q \succ 0$: $\nabla f = Qx - b$, $\mu = \lambda_{\min}(Q)$, $L = \lambda_{\max}(Q)$. \textbf{Condition number} $\kappa = L/\mu \ge 1$; a large $\kappa$ is a long narrow valley.
  \item \textbf{PL condition} (Polyak–Łojasiewicz): $\frac{1}{2}\|\nabla f(x)\|^2 \ge \mu\big(f(x) - f^*\big)$ for all $x$. It follows from $\mu$-strong convexity and makes every stationary point a global minimizer; convexity is not needed ($x^2 + 3\sin^2 x$) nor sufficient ($x^4$).
\end{itemize}

% ---- gd
\cssec{Gradient descent}
\begin{itemize}
  \item \textbf{Iteration} $x_{k+1} = x_k - \alpha\nabla f(x_k)$ along the steepest descent direction $-\nabla f(x_k)$; $d$ is a descent direction if $\nabla f(x)^\top d < 0$.
  \item \textbf{Descent lemma} ($L$-smooth $f$): $f(y) \le f(x) + \nabla f(x)^\top (y - x) + \frac L2\|y - x\|^2$. One step: $f(x^+) \le f(x) - \alpha\big(1 - \frac{L\alpha}{2}\big)\|\nabla f(x)\|^2$, a decrease for $0 < \alpha < 2/L$; the \textbf{safe step} $1/L$ gives $f(x^+) \le f(x) - \frac{1}{2L}\|\nabla f(x)\|^2$.
  \item \textbf{Quadratic}, eigen-coordinates: $z_{k,i} = (1 - \alpha\lambda_i)^k z_{0,i}$. Convergence for every $x_0$ $\iff 0 < \alpha < 2/L$. The \textbf{stability boundary} $2/L$ is excluded: there the stiff coordinate flips sign forever.
  \item \textbf{Contraction factor} $\rho(\alpha) = \max\big(|1 - \alpha\mu|, |1 - \alpha L|\big)$: $\|x_k - x^*\| \le \rho^k\|x_0 - x^*\|$, $f(x_k) - f^* \le \rho^{2k}\big(f(x_0) - f^*\big)$, $\rho(1/L) = 1 - 1/\kappa$. \textbf{Optimal step} $\frac{2}{L + \mu}$: $\rho = \frac{\kappa - 1}{\kappa + 1}$.
\end{itemize}

% ---- rates
\cssec{Rates}
\begin{itemize}
  \item \textbf{Convex}, $\alpha = 1/L$: $f(x_k) - f^* \le \frac{L\|x_0 - x^*\|^2}{2k} = O(1/k)$, since the key inequality $f(x_{k+1}) - f^* \le \frac L2\big(\|x_k - x^*\|^2 - \|x_{k+1} - x^*\|^2\big)$ telescopes.
  \item $\mu$-\textbf{strongly convex} or \textbf{PL}, $\alpha = 1/L$: $f(x_k) - f^* \le (1 - \mu/L)^k\big(f(x_0) - f^*\big)$, from $f(x_{k+1}) \le f(x_k) - \frac{1}{2L}\|\nabla f(x_k)\|^2$ and PL; convexity is not used.
  \item \textbf{Iterations} to cut the gap by the factor $\varepsilon$: $k \ge \frac{\ln\varepsilon}{\ln(1 - 1/\kappa)}$, about $\kappa\ln(1/\varepsilon)$. For $\kappa = 1000$, $\varepsilon = 10^{-6}$: $13809$ at $1/L$, $3454$ at $\frac{2}{L + \mu}$ on a quadratic. A bound, not a forecast: runs are often faster.
\end{itemize}

% ---- armijo
\cssec{Armijo backtracking}
\begin{itemize}
  \item \textbf{Armijo condition}: $f(x + \alpha d) \le f(x) + c\alpha\nabla f(x)^\top d$; for $d = -\nabla f$ the right-hand side is $f(x) - c\alpha\|\nabla f(x)\|^2$. \textbf{Backtracking}: $\alpha = \alpha_0, \beta\alpha_0, \beta^2\alpha_0, \dots$ until it holds; typically $\beta = \tfrac{1}{2}$, $c = 10^{-4}$.
  \item \textbf{Termination} ($L$-smooth $f$): the condition holds for all $\alpha \le 2(1 - c)/L$, so the loop stops with $\alpha \ge \min\big(\alpha_0, 2\beta(1 - c)/L\big)$. Check $\nabla f(x)^\top d < 0$ and cap the reductions.
\end{itemize}

% ---- stopping
\cssec{Stopping rules}
\begin{itemize}
  \item \textbf{Scaled tests}: $\|\nabla f(x_k)\| \le \varepsilon_g\max\big(1, |f(x_k)|\big)$, $\|x_{k+1} - x_k\| \le \varepsilon_x\max\big(1, \|x_k\|\big)$, $|f(x_{k+1}) - f(x_k)| \le \varepsilon_f\max\big(1, |f(x_k)|\big)$; in code $\text{atol} + \text{rtol}\cdot\max(1, s)$.
  \item The change tests shrink with the step, $\|x_{k+1} - x_k\| = \alpha_k\|\nabla f(x_k)\|$, so a tiny step fires them anywhere. \textbf{Combined rule}: success only on the gradient test; a change test alone signals stagnation; cap the iterations and stop when $f$ or $\nabla f$ is not finite.
\end{itemize}

% ---- ad
\cssec{Automatic differentiation}
\begin{itemize}
  \item \textbf{Computational graph}: intermediate variables in topological order; $z = \sin(xy) + x^2$: $u = xy$, $v = \sin u$, $w = x^2$, $z = v + w$. The \textbf{forward pass} computes and stores the values.
  \item \textbf{Reverse pass}: the adjoint $\bar v = \partial z/\partial v$; $\bar z = 1$, the others $0$; visit the nodes in reverse topological order and for each edge $u \to v$ do $\bar u \mathrel{+}= \bar v\,\partial v/\partial u$. Here $\bar x = y\cos u + 2x$, $\bar y = x\cos u$.
  \item \textbf{Accumulation}: a node feeding several others collects one term per edge, $\bar v_i = \sum_{j:\, i \to j} \bar v_j\,\partial v_j/\partial v_i$; $z = x^2 + x^3$ gives $\bar x = 2x + 3x^2$. Assignment keeps only the last term.
  \item \textbf{Forward mode} gives the Jacobian–vector product $J\dot x$ (JVP): $n$ sweeps for a gradient in $\mathbb R^n$. \textbf{Reverse mode} gives the vector–Jacobian product $\bar z^\top J$ (VJP): the whole gradient of a scalar in one sweep. A sweep costs a small multiple of one evaluation of $f$; reverse mode stores every intermediate value (the tape).
\end{itemize}

% ---- fd
\cssec{Finite differences}
\begin{itemize}
  \item \textbf{Central difference} $\frac{f(x + he_i) - f(x - he_i)}{2h}$, error $\approx \frac{h^2}{6}|\partial_i^3 f| + \frac{\varepsilon|f|}{h}$: truncation plus round-off, $\varepsilon$ the machine epsilon. Best $h \sim \varepsilon^{1/3} \approx 6 \cdot 10^{-6}$, error $\sim \varepsilon^{2/3}$. A check of a gradient, not a way to compute it.
\end{itemize}

% ---- logistic
\cssec{Logistic loss}
\begin{itemize}
  \item \textbf{Derivative}: $\frac{d\ell}{ds} = \sigma(s) - y$ for $\ell = \log(1 + e^s) - ys$, $s = w^\top x + b$, $y \in \{0, 1\}$, $\sigma(s) = \frac{1}{1 + e^{-s}}$; so $\nabla_w\ell = \big(\sigma(s) - y\big)x$ and $\frac{\partial\ell}{\partial b} = \sigma(s) - y$.
  \item \textbf{Smoothness} of the mean loss over the rows of $X$: $\nabla^2 f = \frac1n X^\top\operatorname{diag}\big(\sigma_i(1 - \sigma_i)\big)X$ and $\sigma(1 - \sigma) \le \frac{1}{4}$, so $f$ is convex with $L \le \frac{\lambda_{\max}(X^\top X)}{4n}$, attained at $w = 0$.
\end{itemize}

% ---- traps
\cssec{Top traps}
\begin{itemize}
  \item At $\alpha = 2/L$ convergence already fails: it is a boundary, not a step to use; the limit comes from $L = \lambda_{\max}$, and $\mu$ only sets the speed.
  \item Armijo with $+c\alpha\|\nabla f\|^2$ accepts steps that increase $f$; a small change in $x$ or $f$ can mean a tiny step, not a solution.
  \item Reverse pass: $=$ instead of $\mathrel{+}=$ loses terms; $\sigma(s) - y$ is $d\ell/ds$, and the gradient in $w$ carries the factor $x$.
\end{itemize}

\end{multicols}
\end{document}
